\subsection{Elimination, random batching, and algebraic composition soundness}\label{sec:packing:soundness}

The following analysis fixes the hidden witness and
all algebraic values that enter the outgoing equation
\emph{before} the first random batching challenge.
This is necessary to apply Schwartz--Zippel:
if the prover could change the witness or the
masked vector after seeing $\theta$, it could
manufacture a root of the discrepancy polynomial.

\begin{definition}[The three public-value constraints]\label{def:pack:interfaces}
For public $w$-dependent packed values
$(\widehat w_j)_j$, public masked messages
$m\in L^I$, $y'\in L$, a public claimed
evaluation $y_{\mathrm{pub}}$, and hidden
$p,p_{\mathrm{in}},c$, consider:
\begin{align}
 y'&=y_{\mathrm{pub}}+p_{\mathrm{in}},
   \tag{$\mathsf{CLAIM}$}\label{eq:pack:claim}\\
 \sum_vq_vm_v+y'&=\sum_vq_vp_v+p_{\mathrm{in}},
   \tag{$\mathsf{IN}$}\label{eq:pack:in-new}\\
 \sum_ja_j(\theta)\widehat w_j
   &=B_m(\theta)+D_c(\theta).
   \tag{$\mathsf{OUT}$}\label{eq:pack:out-new}
\end{align}
The first two are affine relations in hidden
extension-field values; the third is affine
in hidden packed witness elements and $c$
after $m,\theta$ are public. The further
constraint $c=\mathcal T(p)$ is not
$L$-affine; it will be checked by the
characteristic-two algebraic checker from
Chapter~6.
\end{definition}

\begin{definition}[Partial-evaluation and correction errors]\label{def:pack:errors}
For arbitrary $m,p,c$, define
\[
 e:=m+s^*+p\in L^I,\qquad
 \kappa:=c+c^p\in L^I.
\]
Because the characteristic is two,
$e=0$ is equivalent to $m=s^*+p$,
and $\kappa=0$ to $c=c^p$.
\end{definition}

\begin{lemma}[Exact outgoing discrepancy decomposition]\label{lem:pack:out-error}
For every $\theta\in L^\ell$,
\[
 \boxed{
 B_m(\theta)+D_c(\theta)
   +\sum_ja_j(\theta)\widehat w_j
 = B_e(\theta)+D_\kappa(\theta).
 }
\]
\end{lemma}
\begin{proof}
By the correct packed identity
Theorem~\ref{thm:pack:batch-correct},
$\sum_ja_j\widehat w_j=B_{s^*}$.
By the definition $e=m+s^*+p$,
we have $m=e+s^*+p$. Since $B$
is additive,
$B_m=B_e+B_{s^*}+B_p$.
Substitution into the left side cancels
the two copies of $B_{s^*}$.
The remaining expression is
$B_e+B_p+D_c$.
By Definition~\ref{def:pack:pad-coeff},
$B_p=D_{c^p}$, and linearity of
$D$ in its coefficient table gives
$D_{c^p}+D_c=D_{c^p+c}=D_\kappa$.
\end{proof}

\begin{lemma}[A correct correction isolates the partial-evaluation error]\label{lem:pack:out-isolates}
If $c=c^p$, then the outgoing
equation~\eqref{eq:pack:out-new} holds
exactly when $B_e(\theta)=0$.
\end{lemma}
\begin{proof}
The assumption gives $\kappa=0$,
so $D_\kappa=0$ by its defining sum.
Lemma~\ref{lem:pack:out-error}
shows that the sum of the two
sides of the outgoing equation is
$B_e(\theta)$. In characteristic two,
two field elements are equal precisely
when their sum is zero.
\end{proof}

\begin{lemma}[A false claim and the incoming relations force a row error]\label{lem:pack:false-forces-error}
Assume $\mathsf{CLAIM}$ and $\mathsf{IN}$,
and suppose
$y_{\mathrm{pub}}\neq y_w$.
Then $e\neq0$.
\end{lemma}
\begin{proof}
Suppose instead $e=0$.
By Definition~\ref{def:pack:errors},
$m=s^*+p$. The claim and
incoming equations imply
$y_{\mathrm{pub}}=\sum_vq_v(m_v+p_v)$
by Lemma~\ref{lem:pack:incoming-eliminate}.
Substituting $m+p=s^*$ gives
$y_{\mathrm{pub}}=\sum_vq_vs_v^*$,
which equals $y_w$ by
Lemma~\ref{lem:pack:partial-eval}.
This contradicts the false-claim hypothesis.
\end{proof}

\begin{theorem}[False-claim soundness with exact correction]\label{thm:pack:sound-exact}
Fix $w,m,p,c,y',p_{\mathrm{in}},
y_{\mathrm{pub}},r_{\mathrm{lo}},r_{\mathrm{hi}}$
before uniform $\theta\in L^\ell$.
Suppose $c=c^p$ and
$y_{\mathrm{pub}}\neq y_w$.
Then
\[
 \Pr_\theta[
 \mathsf{CLAIM}\land\mathsf{IN}
               \land\mathsf{OUT}]
 \leq \min\{1,\ell/q\}.
\]
\end{theorem}
\begin{proof}
If either the claim or incoming equation
fails, the acceptance event is empty.
Otherwise Lemma~\ref{lem:pack:false-forces-error}
implies that the fixed vector
$e=m+s^*+p$ is nonzero.
Lemma~\ref{lem:pack:out-isolates}
makes the outgoing equation equivalent
to $B_e(\theta)=0$.
Because $e$ is fixed before $\theta$,
Theorem~\ref{thm:pack:batch-sound}
bounds that event by $\ell/q$.
\end{proof}

\begin{definition}[Combined algebraic acceptance event]\label{def:pack:combined}
Let $\mathsf{Check}_{6}(\theta,z)$ denote
acceptance of the algebraic consistency
checker in Definition~\ref{def:binary:checks},
instantiated with the same hidden pad
vector $p$, candidate coefficient table $c$,
and public challenge $\theta$.
It uses auxiliary hidden vectors/scalars
$t,\rho,\sigma$ and published values $U,A$
subject to the timing in Chapter~6.
Define
\[
 \mathsf{Accept}_{7}:=
 \mathsf{CLAIM}\land\mathsf{IN}
 \land\mathsf{OUT}\land
 \mathsf{Check}_{6}(\theta,z).
\]
\end{definition}

\begin{theorem}[Complete algebraic soundness of packed Ring Switching]\label{thm:pack:algebraic-sound}
Fix $w$, its packed column elements
$(\widehat w_j)$, $m,p,c$,
$y_{\mathrm{pub}},y',p_{\mathrm{in}}$,
and the evaluation points
$r_{\mathrm{lo}},r_{\mathrm{hi}}$
before an independent uniform
$\theta\in L^\ell$.
After $\theta$, an adversary may choose
the auxiliary hidden values and public
trace-zero vector of the Chapter~6 checker,
but all those choices are fixed before
an independent uniform $z\in L$;
it may choose the final public scalar
$A$ after seeing $z$.
Suppose $y_{\mathrm{pub}}\neq y_w$.
Then
\[
 \boxed{
 \Pr_{\theta,z}[\mathsf{Accept}_{7}]
 \leq
 \min\left\{1,\frac{\ell+n}{q}\right\},
 \qquad n=2^\ell,\ q=2^n.}
\]
The bound is for the \emph{algebraic constraints},
not for a concrete externally instantiated
commitment or affine-proof system.
\end{theorem}
\begin{proof}
Because $p,c$ are fixed before
$\theta$, precisely one of the two
cases below holds before challenges
are drawn.

If $c\neq c^p$, acceptance of
$\mathsf{Accept}_{7}$ entails
acceptance of the Chapter~6 checker.
Theorem~\ref{thm:binary:two-challenge}
bounds this probability by
$\min\{1,(\ell+n)/q\}$
under the stated timing, regardless
of the additional incoming and outgoing
constraints.

If $c=c^p$, the Chapter~6 checker
can only reduce the probability of
acceptance. Therefore
$\Pr[\mathsf{Accept}_{7}]$ is at most
the probability of satisfying
$\mathsf{CLAIM}$, $\mathsf{IN}$ and
$\mathsf{OUT}$.
By Theorem~\ref{thm:pack:sound-exact},
this is at most $\min\{1,\ell/q\}$.

The maximum of the bounds of the two
mutually exclusive cases is
$\min\{1,(\ell+n)/q\}$.
This proves the theorem.
\end{proof}

\begin{theorem}[Complete algebraic correctness of packed Ring Switching]\label{thm:pack:algebraic-complete}
Fix a Boolean witness $w$,
its exact packed columns $\widehat w_j$,
and the true public claim
$y_{\mathrm{pub}}=y_w$.
Choose any $p,p_{\mathrm{in}}\in L$
of the declared vector/scalar types,
set $m=s^*+p$, $c=c^p$ and
$y'=y_{\mathrm{pub}}+p_{\mathrm{in}}$.
Choose the Chapter~6 auxiliary
secrets with its honest square relation
and the public messages as in
Theorem~\ref{thm:binary:checker-complete}.
Then $\mathsf{Accept}_{7}$ holds
for every $\theta,z$.
\end{theorem}
\begin{proof}
The honest masked interfaces satisfy
$\mathsf{IN}$ and $\mathsf{OUT}$ by
Theorem~\ref{thm:pack:honest-interfaces};
the defining choice of $y'$
gives $\mathsf{CLAIM}$.
Theorem~\ref{thm:binary:checker-complete}
proves that the Chapter~6 checker
accepts every challenge for the honest
choices of its auxiliary secrets and
public messages. Thus every conjunct
in $\mathsf{Accept}_{7}$ is true.
\end{proof}

\begin{remark}[What the bound excludes]
The theorem assumes exact enforcement of
the claimed field-affine relations,
and that the packed columns, witnesses,
pads, and correction table cannot be
changed after seeing $\theta$.
In an actual system, a commitment or
auxiliary proof protocol must bind
those values and prove each affine
relation with its own soundness bound.
Our theorem does not claim that
the resulting IOP has a
zero-knowledge simulator, nor does it
prove the odd-characteristic
KBFold proximity theorem.
\end{remark}
